% 分野別類題: 波動方程式（ダランベールの解） 解答例
% コンパイル: TEXINPUTS="../../../00_common:$TEXINPUTS" latexmk -xelatex answers.tex
\documentclass[a4paper,11pt]{article}
\usepackage{preamble}
\usepackage{shoakubox}

\begin{document}

\kamokumei{類題演習 解答例 --- 波動方程式（ダランベールの解）}

\daimon{【解法】波動方程式の解き方とダランベールの公式の導出}

\noindent
1次元波動方程式
\[
\frac{\partial^{2} u}{\partial t^{2}} = c^{2}\frac{\partial^{2} u}{\partial x^{2}} \qquad (-\infty < x < \infty,\ t > 0,\ c > 0)
\]
を考える。特性変数
\[
\xi = x - ct, \qquad \eta = x + ct
\]
を導入する（このとき $x = \dfrac{\xi+\eta}{2}$, $t = \dfrac{\eta-\xi}{2c}$）。連鎖律より
\[
\frac{\partial}{\partial x} = \frac{\partial \xi}{\partial x}\frac{\partial}{\partial \xi} + \frac{\partial \eta}{\partial x}\frac{\partial}{\partial \eta}
= \frac{\partial}{\partial \xi} + \frac{\partial}{\partial \eta},
\qquad
\frac{\partial}{\partial t} = \frac{\partial \xi}{\partial t}\frac{\partial}{\partial \xi} + \frac{\partial \eta}{\partial t}\frac{\partial}{\partial \eta}
= -c\frac{\partial}{\partial \xi} + c\frac{\partial}{\partial \eta}
\]
したがって
\[
\frac{\partial^{2}}{\partial x^{2}} = \left(\frac{\partial}{\partial \xi} + \frac{\partial}{\partial \eta}\right)^{2}
= \frac{\partial^{2}}{\partial \xi^{2}} + 2\frac{\partial^{2}}{\partial \xi\,\partial \eta} + \frac{\partial^{2}}{\partial \eta^{2}}
\]
\[
\frac{\partial^{2}}{\partial t^{2}} = \left(-c\frac{\partial}{\partial \xi} + c\frac{\partial}{\partial \eta}\right)^{2}
= c^{2}\left(\frac{\partial^{2}}{\partial \xi^{2}} - 2\frac{\partial^{2}}{\partial \xi\,\partial \eta} + \frac{\partial^{2}}{\partial \eta^{2}}\right)
\]
これを波動方程式 $u_{tt} = c^{2}u_{xx}$ に代入すると、$\xi^{2}$, $\eta^{2}$ の項は両辺で打ち消し合い
\[
c^{2}\left(-4\frac{\partial^{2} u}{\partial \xi\,\partial \eta}\right) = 0
\quad\Longleftrightarrow\quad
\frac{\partial^{2} u}{\partial \xi\,\partial \eta} = 0
\]
となる。この式は
\[
\frac{\partial}{\partial \eta}\left(\frac{\partial u}{\partial \xi}\right) = 0
\]
と読める。すなわち「$\dfrac{\partial u}{\partial \xi}$ を $\eta$ で微分するとゼロ」だから、$\dfrac{\partial u}{\partial \xi}$ は $\eta$ に依存しない。つまり $\xi$ のみの関数であり、これを $F(\xi)$ とおく:
\[
\frac{\partial u}{\partial \xi} = F(\xi) \qquad (F \text{ は任意関数})
\]
次にこれを、$\eta$ を固定したまま $\xi$ で積分する。偏微分の積分なので「積分定数」は $\xi$ に関して定数であればよく、$\eta$ の任意の関数 $g(\eta)$ になりうることに注意すると
\[
u(\xi,\eta) = \underbrace{\int F(\xi)\,d\xi}_{=\,f(\xi) \text{ とおく}} + g(\eta)
\]
$F$ が任意関数なのでその原始関数 $f$ も任意関数である。よって
\[
u(\xi,\eta) = f(\xi) + g(\eta)
\]
（$f$, $g$ は任意の $C^{2}$ 級関数）を得る。元の変数に戻せば、一般解は
\[
u(x,t) = f(x-ct) + g(x+ct)
\]
となる。これは、右向きに速さ $c$ で進む波 $f(x-ct)$ と左向きに速さ $c$ で進む波 $g(x+ct)$ の重ね合わせを表す。

初期条件 $u(x,0) = \varphi(x)$, $u_{t}(x,0) = \psi(x)$ が与えられたとき、$f$, $g$ を定める。
\[
u(x,0) = f(x) + g(x) = \varphi(x)
\]
\[
u_{t}(x,t) = -c f'(x-ct) + c g'(x+ct)
\quad\Longrightarrow\quad
u_{t}(x,0) = -c f'(x) + c g'(x) = \psi(x)
\]
第2式を $x$ で積分すると（積分定数を $A$ とする）
\[
-f(x) + g(x) = \frac{1}{c}\int_{0}^{x} \psi(s)\,ds + A
\]
これと $f(x)+g(x)=\varphi(x)$ を連立して解くと
\[
f(x) = \frac{1}{2}\varphi(x) - \frac{1}{2c}\int_{0}^{x}\psi(s)\,ds - \frac{A}{2},
\qquad
g(x) = \frac{1}{2}\varphi(x) + \frac{1}{2c}\int_{0}^{x}\psi(s)\,ds + \frac{A}{2}
\]
これらを $u(x,t) = f(x-ct) + g(x+ct)$ に代入し、定数 $A$ が消えることを確かめると
\[
u(x,t) = \frac{\varphi(x-ct) + \varphi(x+ct)}{2} + \frac{1}{2c}\int_{x-ct}^{x+ct} \psi(s)\,ds
\]
を得る。これが\textbf{ダランベールの公式}である。

\clearpage
\begin{mondaiboxS}{問1}
波動方程式 $u_{tt} = c^{2}u_{xx}$ の一般解が $u(x,t) = f(x-ct) + g(x+ct)$ の形で表されることを、特性変数 $\xi = x-ct$, $\eta = x+ct$ を用いて示せ。
\end{mondaiboxS}
\kaitou
$\xi = x-ct$, $\eta = x+ct$ とおくと、連鎖律より
\[
\frac{\partial}{\partial x} = \frac{\partial}{\partial \xi} + \frac{\partial}{\partial \eta}, \qquad
\frac{\partial}{\partial t} = -c\frac{\partial}{\partial \xi} + c\frac{\partial}{\partial \eta}
\]
であるから
\[
\frac{\partial^{2} u}{\partial x^{2}} = u_{\xi\xi} + 2u_{\xi\eta} + u_{\eta\eta}, \qquad
\frac{\partial^{2} u}{\partial t^{2}} = c^{2}\left(u_{\xi\xi} - 2u_{\xi\eta} + u_{\eta\eta}\right)
\]
これらを $u_{tt} = c^{2}u_{xx}$ に代入すると
\[
c^{2}\left(u_{\xi\xi} - 2u_{\xi\eta} + u_{\eta\eta}\right) = c^{2}\left(u_{\xi\xi} + 2u_{\xi\eta} + u_{\eta\eta}\right)
\quad\Longrightarrow\quad
-4c^{2}u_{\xi\eta} = 0
\quad\Longrightarrow\quad
u_{\xi\eta} = 0
\]
（$c \neq 0$ より）。$u_{\xi\eta} = \dfrac{\partial}{\partial \eta}\left(\dfrac{\partial u}{\partial \xi}\right) = 0$ を $\eta$ について積分すると、$\dfrac{\partial u}{\partial \xi}$ は $\eta$ に依存しない、すなわち $\xi$ のみの関数 $F(\xi)$ に等しい。さらに $\xi$ について積分すると
\[
u(\xi,\eta) = \int F(\xi)\,d\xi + g(\eta) =: f(\xi) + g(\eta)
\]
（$f$, $g$ は任意の $C^{2}$ 級関数）。元の変数 $x, t$ に戻すと
\[
\boxed{\; u(x,t) = f(x-ct) + g(x+ct) \;}
\]
（検算: $u_x = f' + g'$, $u_{xx} = f'' + g''$。$u_t = -cf' + cg'$, $u_{tt} = c^{2}f'' + c^{2}g'' = c^{2}u_{xx}$ となり、確かに波動方程式を満たす。）

\clearpage
\begin{mondaiboxS}{問2}
波動方程式 $u_{tt} = c^{2}u_{xx}$ が $u(x,0) = \varphi(x)$, $u_{t}(x,0) = 0$ を満たすとき $u(x,t)$ を求めよ。特に $\varphi(x) = e^{-x^{2}}$, $c=2$ の場合を具体的に求めよ。
\end{mondaiboxS}
\kaitou
ダランベールの公式
\[
u(x,t) = \frac{\varphi(x-ct) + \varphi(x+ct)}{2} + \frac{1}{2c}\int_{x-ct}^{x+ct} \psi(s)\,ds
\]
において $\psi \equiv 0$ とすると積分項が消え
\[
\boxed{\; u(x,t) = \frac{\varphi(x-ct) + \varphi(x+ct)}{2} \;}
\]
$\varphi(x) = e^{-x^{2}}$, $c = 2$ を代入すると
\[
\boxed{\; u(x,t) = \frac{1}{2}\left(e^{-(x-2t)^{2}} + e^{-(x+2t)^{2}}\right) \;}
\]
（検算: $t=0$ のとき $u(x,0) = \dfrac{1}{2}(e^{-x^2}+e^{-x^2}) = e^{-x^2} = \varphi(x)$ OK。
$u_t = \dfrac{1}{2}\left(4(x-2t)e^{-(x-2t)^2} - 4(x+2t)e^{-(x+2t)^2}\right)$ より $u_t(x,0) = \dfrac{1}{2}(4xe^{-x^2}-4xe^{-x^2}) = 0$ OK。
また各項は $x\mp 2t$ のみの関数なので問1の形 $f(x-ct)+g(x+ct)$ に含まれ、自動的に波動方程式 $u_{tt}=4u_{xx}$ を満たす。）

\clearpage
\begin{mondaiboxS}{問3}
波動方程式 $u_{tt} = c^{2}u_{xx}$ が $u(x,0) = 0$, $u_{t}(x,0) = \psi(x)$ を満たすとき $u(x,t)$ を求めよ。特に $\psi(x) = \cos x$, $c=3$ の場合を具体的に求めよ。
\end{mondaiboxS}
\kaitou
ダランベールの公式において $\varphi \equiv 0$ とすると
\[
\boxed{\; u(x,t) = \frac{1}{2c}\int_{x-ct}^{x+ct} \psi(s)\,ds \;}
\]
$\psi(x) = \cos x$, $c = 3$ を代入すると
\[
u(x,t) = \frac{1}{6}\int_{x-3t}^{x+3t} \cos s\,ds
= \frac{1}{6}\Big[\sin s\Big]_{x-3t}^{x+3t}
= \frac{1}{6}\left(\sin(x+3t) - \sin(x-3t)\right)
\]
和積公式 $\sin A - \sin B = 2\cos\left(\dfrac{A+B}{2}\right)\sin\left(\dfrac{A-B}{2}\right)$ より $\sin(x+3t)-\sin(x-3t) = 2\cos x \sin 3t$ であるから
\[
\boxed{\; u(x,t) = \frac{1}{3}\cos x \sin 3t \;}
\]
（検算: $u(x,0) = \dfrac{1}{3}\cos x \sin 0 = 0$ OK。
$u_t = \cos x \cos 3t$ より $u_t(x,0) = \cos x = \psi(x)$ OK。
$u_{tt} = -3\cos x \sin 3t$。$u_x = -\dfrac{1}{3}\sin x \sin 3t$, $u_{xx} = -\dfrac{1}{3}\cos x \sin 3t$ より $c^{2}u_{xx} = 9 \times \left(-\dfrac{1}{3}\cos x \sin 3t\right) = -3\cos x \sin 3t = u_{tt}$ OK。波動方程式を満たす。）

\end{document}
